\subsection{Тензорезисторные датчики}
Их работа основана на тензоэффекте: изменение электрического сопротивления тензорезистора при его деформации. 

\begin{equation}
	K = \frac{\sfrac{\Delta R}{R}}{\epsilon}
\end{equation}
- коэффициент относительной тензочувствительности.

\begin{equation}
	\frac{\Delta R}{R} = \frac{\Delta \rho}{\rho} + \frac{\Delta l}{l} - \frac{\Delta S}{S} = \frac{\Delta \rho}{\rho} + \frac{\Delta l}{l} - \frac{2\Delta r}{r} = \frac{\Delta \rho}{\rho} + (1 + 2\mu)\epsilon
\end{equation}

где $\mu$ - коэффициент Пуассона. Первое слагаемое - физическая тензочувствиельность, а второе - геометрическая тензочувствительность. 

Для металлических проводников 

\begin{equation}
	\Delta \rho = \rho [\pi_{11} \sigma_{1} + \pi_{12}(\sigma_{2} + \sigma_{3})]
\end{equation}
где $\sigma_{123}$ - компоненты нормальных напряжений, $\pi_{11} \ pi_{12}$ - тензорезистивные коэффициенты.

При линейных напряжениях $\sigma_{1} = \sigma, \ \sigma_{2} = \sigma_{3} = 0$.
Откуда

\begin{equation}
	\frac{\Delta \rho}{\rho} = \pi_{11} \sigma = \pi_{11}E\epsilon    
\end{equation}
где Е - модуль Юнга.

\newpage